% Auto-generated by PIO-Net Task 4.2 for ACM SenSys 2027
% Formatted for ACM Double-Column Format (acmart sigconf)

% --- Figure 1: 12-Class Confusion Matrix ---
\begin{figure}[t]
  \centering
  \includegraphics[width=\columnwidth]{figures/pionet_confusion_matrix.pdf}
  \caption{Normalized 12-class confusion matrix for Stage-2 INT8 \textsc{DSCNNMicro} on the MLPerf Tiny Speech Commands dataset (Macro Recall / Balanced Accuracy: 90.8\%). Main diagonal cells depict class-specific recall; values are rounded to the nearest integer. The orange dashed line isolates the 10 target keywords from the background rejection classes (\texttt{\_silence\_} and \texttt{\_unknown\_}).}
  \label{fig:confusion_matrix}
\end{figure}

% --- Figure 2: Detection Error Tradeoff (DET) Curves ---
\begin{figure}[t]
  \centering
  \includegraphics[width=\columnwidth]{figures/pionet_det_curves.pdf}
  \caption{Detection Error Tradeoff (DET) curves across three acoustic noise categories. The primary bottom axis displays calibrated False Wakes per Hour (FW/hr), while the top axis shows the corresponding False Alarm Probability $P_{\mathrm{FA}}$ per 1.0-second sliding inference window. Long-term continuous rates are parametrically projected from empirical score distributions. Stationary ambient noise achieves lower false alarm rates due to the absence of resonant human vocal formants present in office babble and mechanical harmonics. Star glyphs highlight the selected operating threshold ($\theta = 22$), achieving $\le 1.0\ \mathrm{FW/hr}$ with $< 20\%$ Miss Probability ($> 80\%$ Recall).}
  \label{fig:det_curves}
\end{figure}

% --- Figure 3: Cascaded Energy Waterfall & Battery Lifetime ---
\begin{figure}[t]
  \centering
  \includegraphics[width=\columnwidth]{figures/pionet_energy_waterfall.pdf}
  \caption{Cascaded system power breakdown and projected operating lifetime on standard portable batteries. (a) PIO-Net achieves a 72.4\% power reduction over continuous CPU execution by keeping the RP2040 ARM Cortex-M0+ clock-gated in deep sleep during quiescent monitoring. (b) Operating lifetimes are projected from usable electrochemical cell capacity down to the RP2040 supply cutoff (170 mAh for CR2032, 450 mAh for LiPo 500, 850 mAh for 2x AAA) at nominal cell voltages.}
  \label{fig:energy_waterfall}
\end{figure}

% --- Table 1: End-to-End System Performance Benchmark ---
\begin{table}[t]
  \centering
  \small
  \caption{End-to-End Wake-Up and Energy Performance on RP2040}
  \label{tab:system_comparison}
  \begin{tabular}{lcccc}
    \toprule
    \textbf{Pipeline Configuration} & \textbf{Duty Cycle} & \textbf{Power} & \textbf{2x AAA Life} & \textbf{FWR} \\
    \midrule
    Config A (Continuous CPU) & 100.0\% & 65.01 mW & 39.2 h & 0.0 /h \\
    Config B (Periodic 20 ms)  &  12.6\% & 26.67 mW & 95.6 h & 0.0 /h \\
    \textbf{PIO-Net (Cascaded)} & \textbf{0.0\%}$^\dagger$ & \textbf{17.95 mW} & \textbf{142.1 h} & \textbf{0.0 /h} \\
    \bottomrule
    \multicolumn{5}{l}{\footnotesize $^\dagger$ Quiescent CPU duty cycle during background monitoring.}
  \end{tabular}
\end{table}
